A spacecraft is launched from the Earth and it is quickly accelerated to its
maximum velocity in the direction of the heliocentric orbit of the Earth, such
that its orbit is a parabola with the Sun at its focus point, and grazes the
Earth orbit. Take the orbit of the Earth and Mars as circles on the same plane,
with radius of $r_E = 1$\,AU and $r_M = 1.5$\,AU, respectively. Make the
following approximation: during most of the flight only the gravity from the
Sun needs to be considered. The figure below shows the trajectory of the
spacecraft (not in scale); the inner circle is the orbit of the Earth, the
outer circle is the orbit of Mars.

\ioaafig{2010_TH_L16-f1.pdf}

\begin{description}
\item[(a)] What is the angle $\psi$ between the path of the spacecraft and the
  orbit of the Mars (see the figure) as it crosses the orbit of the Mars,
  without considering the gravity effect of the Mars?
\item[(b)] Suppose the Mars happens to be very close to the crossing point at
  the time of the crossing, from the point of view of an observer on Mars, what
  is the approaching velocity and direction of approach (with respect to the
  Sun) of the spacecraft before it is significantly affected by the gravity of
  the Mars?
\end{description}
